% preamble.tex — paquetes y comandos globales del informe 3DM
\usepackage[spanish]{babel}
\usepackage[utf8]{inputenc}
\usepackage[T1]{fontenc}
\usepackage{lmodern}

% Matemáticas
\usepackage{amsmath, amssymb, amsthm}

% Tipografía de tablas y figuras
\usepackage{booktabs}
\usepackage{graphicx}
\usepackage{float}

% Hipervínculos
\usepackage[hidelinks]{hyperref}

% Algoritmos
\usepackage[ruled,vlined,linesnumbered]{algorithm2e}
\SetKwInput{KwIn}{Entrada}
\SetKwInput{KwOut}{Salida}

% Código fuente (para snippets inline)
\usepackage{listings}
\usepackage{xcolor}

% Geometría de página
\usepackage[margin=2.5cm]{geometry}

% Bibliografía con natbib (compatible con tectonic sin biber)
\usepackage[numbers,sort&compress]{natbib}

% Entornos matemáticos
\newtheorem{lemma}{Lema}[section]
\newtheorem{theorem}[lemma]{Teorema}
\newtheorem{corollary}[lemma]{Corolario}
\theoremstyle{definition}
\newtheorem{definition}[lemma]{Definición}
\theoremstyle{remark}
\newtheorem*{remark}{Nota}

% Operadores matemáticos
\DeclareMathOperator*{\argmax}{arg\,max}
\DeclareMathOperator{\OPT}{OPT}

% Abreviaturas útiles
\newcommand{\threedm}{3\nobreakdash-Dimensional Matching}
\newcommand{\PNPO}{\ensuremath{\text{P} \neq \text{NP}}}

% Estilo de código
\lstset{
  basicstyle=\ttfamily\small,
  keywordstyle=\color{blue},
  commentstyle=\color{gray},
  breaklines=true,
  frame=single,
  numbers=left,
  numberstyle=\tiny\color{gray},
}
